Just-in-time software defect prediction (JIT-SDP) labels almost
universally descend from the SZZ algorithm, which infers
defect-introducing commits by blaming the lines a bug fix modifies. That
inference fails whenever bug-fixing commits are tangled, and the field
has no measurement of what the resulting noise costs a model under
realistic deployment conditions --- in part because models are
conventionally evaluated against the same SZZ labels that trained them.

This thesis separates four quantities across 21 Apache Java projects and
27,319 commits, holding out an independently constructed reference label
set so that a model trained on SZZ can be scored against something other
than SZZ.

\textbf{Label quality.} Six SZZ variants are scored over an identical
denominator. No variant exceeds \textbf{27.2\% precision}. Noise is
class-conditional rather than symmetric (false-alarm rates 6.7--26.3\%,
miss rates 35.9--73.3\%). Variants agree with one another at \ensuremath{\kappa} up to
\textbf{0.933} while agreeing with the reference at \ensuremath{\kappa}
\textbf{0.158--0.202} --- they share failure modes, so mutual agreement
is not evidence of correctness.

\textbf{Evaluation protocol.} Random \emph{k}-fold cross-validation
inflates a high-capacity model by \textbf{+0.137 MCC} (21 of 21
projects); scoring against the training labels inflates it by a further
\textbf{+0.230}. The two protocol effects together are roughly nine
times the label-source effect, and both scale with model capacity. A
result reading 0.41 MCC under the field's standard configuration is
\textbf{0.10} under a defensible one.

\textbf{Mechanism.} A controlled repair experiment shows that removing
SZZ's false positives recovers \textbf{+0.040 MCC} while restoring its
false negatives is statistically \emph{equivalent to no repair}. An
error-matched control establishes this as an effect of error
\textbf{volume}, not per-label severity. The mechanism is measured: the
online learner's oversampling rate is a near-perfect inverse of
delivered defect-label supply (\ensuremath{\rho} = \ensuremath{-}1.000), so the machinery that
corrects class imbalance is the machinery that amplifies wrong labels.

\textbf{Mitigation.} Two pre-registered experiments. A train-time
false-positive filter passes its non-degradation gate but no
confirmatory test survives correction; the gap between the
oracle-assisted repair (+0.040) and its deployable approximation
(+0.017) is the price of not knowing which labels are wrong, since a
confidence signal cannot separate a wrong label from an unlearned
pattern. A second registration then confirms that confidence damping
beats both baseline and filter --- and a control shows the entire
benefit is the \textbf{absence of oversampling-rate boosting} rather
than any defence: plain online bagging matches every noise-aware variant
tested and beats the boosted learner in \textbf{18 of 18 projects} under
false-positive-heavy noise.

This is independent confirmation, by intervention, of the amplification
mechanism established by measurement, and it yields the thesis's one
actionable recommendation. \textbf{Imbalance correction and label noise
interact, and the standard remedy for the first makes the second worse:}
under SZZ-derived labels, do not boost the oversampling rate on
prediction bias.

Improving the labels remains worth more than improving the learner ---
but the cheapest effective intervention is neither.
